\section{Related Work}
\label{sec:related}

A growing benchmark literature probes whether LLMs retain
information over long inputs, but does so almost exclusively
through the verbal channel. LongMemEval \citep{Wu2025LongMemEval}
constructs multi-session dialogue probes that score retrieval of
prior facts through question answering; RULER
\citep{Hsieh2024RULER} measures needle-in-a-haystack retrieval as a
function of context length and content type; Lost-in-the-Middle
\citep{Liu2024LostInMiddle} characterises positional bias in
document QA; InfiniteBench \citep{Zhang2024InfiniteBench} extends
the same recall paradigm to 100k-token regimes. These benchmarks
share a common evaluation surface: a textual probe is presented and
the model's textual answer is scored. They do not measure what the
model \emph{does} when given the same context inside an agent loop.
A second line of work probes action grounding inside agent loops:
MemoryArena \citep{He2026MemoryArena} and MemoryAgentBench
\citep{Hu2025MemoryAgentBench} evaluate whether memory-augmented
agents can ground prior context into correct actions on benchmark
tasks, but the episodes are session-discrete---each task is a
fresh prompt or a fresh agent run, and memory is tested as
retrieval-into-action rather than as the sustained joint behaviour
of two channels of the same continuous process. ReAct
\citep{Yao2023ReAct} and Reflexion \citep{Shinn2023Reflexion}
introduced the interleaved reasoning-and-action paradigm that these
benchmarks evaluate, but report aggregate task success rather than
channel-level dissociation. Recent work on diminishing returns in
LLM agents \citep{Sinha2025IllusionDiminishingReturns} characterises
performance ceilings without separating the verbal and action
contributions to those ceilings. Anecdotal reports of
verbal--action dissociation appear in post-mortems of
production-grade agent deployments: the team behind Cicero
\citep{FAIR2022Cicero} documented cases in which dialogue and
action diverged under pressure, and Anthropic's multi-agent system
\citep{AnthropicMultiAgent2025} reports that sub-agents
occasionally produce summaries that overstate or misalign with the
actions they took. These reports are descriptive: they note that
the phenomenon occurs, but offer no controlled probe, no
quantitative threshold, and no testbed on which to study its
horizon-dependence or its tier-dependence. Our contribution sits in
the gap: a joint-axis probe on tick-continuous structured-output
context, with a pre-registered threshold ($|\Delta| > 0.3$ on a
behaviourally discounted dual probe), a reproducible substrate, and
pilot results that characterise the phenomenon's magnitude across
deployment tiers.
